\documentclass[10pt,letterpaper]{article}
\usepackage[margin=0.42in]{geometry}
\usepackage[T1]{fontenc}
\usepackage[utf8]{inputenc}
\usepackage{lmodern}
\usepackage{xcolor}
\usepackage{titlesec}
\usepackage{enumitem}
\usepackage{tabularx}
\usepackage[hidelinks]{hyperref}

% Contact information (kept inline so this revision compiles from one source file)
\newcommand{\mylocation}{Murphy, NC}
\newcommand{\myemail}{consav04@gmail.com}
\newcommand{\myphone}{+1-828-541-9487}
\newcommand{\mylinkedin}{connorsavugot}
\newcommand{\mygithub}{Clem085}

\pagestyle{empty}
\setlength{\parindent}{0pt}
\setlength{\parskip}{0pt}

\titleformat{\section}{\large\bfseries}{\thesection}{1em}{}[\vspace{-0.5em}\rule{\textwidth}{0.7pt}]
\titlespacing{\section}{0pt}{0.5em}{0.25em}

\newlist{compactitem}{itemize}{1}
\setlist[compactitem]{nosep, leftmargin=0.3in, label=\textbullet, itemsep=0.05em}

\begin{document}
\fontsize{8.5}{9.2}\selectfont

\begin{center}
  {\LARGE\textbf{Connor Savugot}}\\[2pt]
  {\small \mylocation \enspace|\enspace \href{mailto:\myemail}{\myemail} \enspace|\enspace \href{tel:\myphone}{\myphone} \\
  \href{https://linkedin.com/in/\mylinkedin}{linkedin.com/in/\mylinkedin} \enspace|\enspace \href{https://github.com/\mygithub}{github.com/\mygithub}}
\end{center}

\section{Technical Skills}
\begin{compactitem}
  \item \textbf{Languages:} C, C++, Python, Bash, MATLAB, VBA
  \item \textbf{Embedded Systems:} FreeRTOS, Embedded Linux, systemd services, IPMI, external watchdogs, BMS integration, Yocto/Arago, device trees, board bring-up
  \item \textbf{Platforms:} TI MSP430, TI TMS320, TI AM62-class / ARM Linux SoCs, STM32, dsPIC
  \item \textbf{Protocols \& Peripherals:} CAN, CAN FD, UART, RS232, SPI, I2C, ADC, PWM, GPIO, timers, JTAG
  \item \textbf{Hardware \& Tools:} Through-hole and SMT soldering, PCB rework, breadboard prototyping, KiCad, Git/GitLab, JTAG, oscilloscopes, logic analyzers
\end{compactitem}

\section{Professional Experience}
\textbf{AEGIS Power Systems} \hfill Current\\
Embedded Firmware Engineer
\begin{compactitem}
  \item Developed firmware control for custom bidirectional buck-boost power supplies using 24 independently configurable PWM channels across transformer-coupled input and output stages.
  \item Built a centralized configuration framework that translated requested input and output voltages into coordinated PWM timing, duty-cycle, and switching behavior across paired converter channels; integrated ADC feedback and safe-operating constraints.
  \item Evaluated BMS, monitoring, interface, and support ICs against electrical limits, host protocols, protection features, footprints, tool support, and sourcing risk; converted datasheet and reference-design findings into firmware and integration plans.
  \item Root-caused I2C, CAN/CAN FD, SPI, and UART faults by correlating register transactions and analyzer captures with schematics, timing requirements, bus configuration, and known-good hardware; implemented focused tests for peripheral-by-peripheral bring-up.
  \item Characterized IPMI 2.0 and chassis-management services on embedded power hardware, validating IPMB addressing plus sensor and FRU access while separating transport, session, authentication, and firmware-compatibility failures.
  \item Created external-watchdog and systemd services, device-tree and Yocto/Arago support, and Python/Bash tooling for reproducible monitoring, flashing, automated validation, and field diagnostics on TI AM62-class Linux systems.
\end{compactitem}

\textbf{AEGIS Power Systems} \hfill May 2024 -- Aug 2024\\
Computer Science Intern
\begin{compactitem}
  \item Built a secured chroot environment and restricted SSH command-line interface on a stripped TI Yocto/Arago system, using custom Bash/Python commands to inspect and control embedded power hardware.
  \item Created Python and Bash firmware flashing utilities with packetized transfer, checksum validation, and SPI reflashing support for TI TMS320 targets.
  \item Reverse-engineered serial command transactions from limited vendor documentation, generated checksums, and helped construct a relay, resistive-load, and power-supply test circuit for burn-in monitoring.
  \item Authored Git and embedded-system guides to standardize firmware delivery, testing, and toolchain onboarding.
\end{compactitem}

\textbf{TEAM Industries} \hfill Jun 2022 -- Aug 2022\\
Engineering Intern
\begin{compactitem}
  \item Automated Excel-based quality-control reporting and revision tracking using VBA macros and structured technical worksheets.
  \item Produced technical documentation that clarified engineering changes and supported manufacturing process reviews.
\end{compactitem}

\section{Selected Technical Projects}
\textbf{EV Active Sensor Adapter} \hfill Senior Design, MSP430FR2355\\
\begin{compactitem}
  \item Defined the end-to-end architecture and interface specifications for an EV active-sensor adapter integrating an MSP430, BMS and power-management hardware, sensors, custom PCBs, embedded firmware, and system telemetry.
  \item Evaluated BMS and supporting components against electrical, sensing, protection, communication, packaging, and sourcing requirements; documented trade studies and translated selected-device requirements into hardware and firmware interfaces.
  \item Prototyped interrupt-driven SPI master/slave register transactions and integrated ADC acquisition, UART diagnostics, timers, GPIO, and LCD telemetry while validating each subsystem before full-system integration.
  \item Maintained requirements, architecture diagrams, interface definitions, test plans, design decisions, and progress reports; coordinated work across team members and communicated specifications, risks, and results during design reviews.
\end{compactitem}

\section{Education}
\textbf{North Carolina State University} \hfill May 2026\\
B.S. Computer Engineering

\section{Technical Leadership}
\textbf{Lead Instructor, Project-Based Electronics Workshop} \hfill Two Semesters\\
\begin{compactitem}
  \item Co-led a project-based class from concept through final demonstration, guiding students as they scoped, designed, debugged, and completed original electronics projects.
  \item Developed and taught through-hole and SMT soldering theory, safety, and hands-on technique using two custom NCSU IEEE flashing-light PCB keepsakes: one through-hole and one surface-mount.
\end{compactitem}

\textbf{Embedded Systems Teaching Assistant} \hfill NC State University\\
\begin{compactitem}
  \item Delivered supplemental lectures and test preparation; taught embedded theory, code debugging, lab techniques, and systematic hardware bring-up while advising students on projects, coursework, and ECE pathways.
  \item Used question-driven, one-on-one coaching to help students diagnose their own problems and form lasting professional relationships, raising the class average from its typical C+ to an A-.
\end{compactitem}

\end{document}
